
\chapter{Algorithmic Performance} 

For algorithms whose performance\footnote{Although this section focuses primarily on running 
time, the same methods of analysis can also be applied to other resources such as memory usage.} 
depends on the input size or other input parameters, we usually estimate the number of 
elementary operations performed as a function of the input size \(n\) or the maximum 
value of input $A_{max}$, usually denoted by $A$. This gives a mathematical way 
to compare algorithms independent of hardware or implementation details.\\


\noindent A classical example is \texttt{BubbleSort}. The algorithm repeatedly traverses the 
array, comparing adjacent elements and swapping them whenever they are in the 
wrong order. After each pass, the largest remaining element moves to its 
correct position near the end of the array. With some thought it can be proved 
that the total number of swaps performed by \texttt{BubbleSort} is exactly 
equal to the number of inversions $N_{inv}$ in the input array. 
(An inversion is a pair of indices \((i,j)\) such that \(i<j\) but \(a_i>a_j\).) By taking 
swaps to be one operation taking constant time, we can estimate runtime 
of \texttt{BubbleSort} to be $T_{run}\propto N_{inv}$

\section{Operations}
In complexity analysis, we must know which operations are to be taken as a unit of time. This is sometimes a 
bit subjective since arithmetic operations (+,-,$\times$,$\div$) can be taken as constant time 
if we operate under the assumption that we are using fixed-length numbers like 32-bit integers.

Regardless, here is a list of the common operations taken as a unit of time in complexity analysis.

\begin{itemize}

\item Arithmetic operations (+,-,$\times$,$\div$) 

\item Comparisons ($<$,$>$,$\leq$,$\geq$)

\item Boolean operations ($\land$,$\lor$,$\neg$)

\item Bitwise operations ($\&$,$\|$,$\neg$)

\item Variable read/write

\item Array read/write 

\end{itemize}


\section{Worst Case Analysis}

In worst case analysis, we determine the maximum possible number of operations over all valid 
inputs of size \(n\). For \texttt{BubbleSort}, the worst case occurs when the array is 
sorted in reverse order, where every pair is inverted. 

\[N_{inv,max}=\binom{n}{2}=
\frac{n(n-1)}{2}
\]

\noindent Worst case analysis is useful because it guarantees an upper bound 
on the running time regardless of the input provided.

\section{Average Case Analysis} 

In average case analysis, we compute the expected number of operations 
assuming that all valid inputs are equally likely. For \texttt{BubbleSort}, this corresponds 
to analysing the expected number of inversions in a random permutation of size \(n\), where all 
$n!$ permutations are equally likely.

\noindent Define the indicator\footnote{Refer \autoref{indicator_variables}} variable $I_{ij}$ 
for each pair \((i,j)\) as
\[
I_{ij}=
\begin{cases}
1,&\text{if }(i,j)\text{ is an inversion}\\
0,&\text{otherwise}
\end{cases}
\]
The point of defining the indicator variable is to notice that the number 
of inversions is the sum of all $I_{ij}$. Formally we can write this as 
\[
N_{\mathrm{inv}}=\sum_{1\le i<j\le n}I_{ij}
\]
For every pair of indices \((i,j)\) with \(i<j\), the probability that the pair 
forms an inversion is $\frac{1}{2}$ since for every permmutation with $a_i>a_j$, 
there is a permuation with $a_j>a_i$ and vice versa.
\[
\Pr\big((i,j)\text{ is an inversion}\big)=\frac{1}{2} \implies \mathbb{E}[I_{ij}]
=
\frac{1}{2}
\]
Using linearity of expectation\footnote{Refer \autoref{linearity_of_expectation}}, 
\[
\mathbb{E}[N_{\mathrm{inv}}]
=
\mathbb{E}\left[\sum_{1\le i<j\le n}I_{ij}\right]
=
\sum_{1\le i<j\le n}\mathbb{E}[I_{ij}]
=
\sum_{1\le i<j\le n}\frac{1}{2}
=
\frac{n(n-1)}{4}
\]

\noindent A valuable strategy which we can deduce from this analysis is to randomise the 
input sequence. In case of an adversary\footnote{An adversary is someone who deliberately 
gives inputs to make your algorithm perform badly. Randomising the input is one possible defense 
against adversaries.}, this can be useful because we 
can expect an improvement in time by factor of $2$.

\section{Time Complexity Analysis}

Let us first define the asymptotic complexity of a function \(f(n)\) in terms of \(n\). My favourite way to 
think about complexity is that we try to find the ``simplest'' function \(g(n)\) such that
\[
\lim_{n\to\infty}\frac{f(n)}{g(n)}=k, \;\;
k\in\mathbb{R}_{>0}
\]

This is compactly represented as $f(n)=\Theta(g(n))$. If via inequalities we 
are able to find $f(n) \le h(n)$ or $f(n) < h(n)$ and $h(n)=\Theta(g(n))$, then we can represent that by $f(n) \le \Theta(g(n))$. 
Similarly we define what it means to say that $f(n) \ge \Theta(g(n))$.\\


\noindent By applying Sandwich theorem on the ratio $f(n)/g(n)$, we can conclude that if $f(n) \ge \Theta(g(n))$ 
and $f(n) \le \Theta(g(n))$ then $f(n)=\Theta(g(n))$.\\


\noindent Since "simplest" function is not really well defined, it is best 
to look at some examples and gain intuition!
\begin{center}
\begin{tabular}{|c|c|}
\hline
Function & Asymptotic Complexity \\
\hline
$7$ & $\Theta(1)$ \\
\hline
$n^2+100n+1$ & $\Theta(n^2)$ \\
\hline
$n+\log n$ & $\Theta(n)$ \\
\hline
$\log(n!)$ & $\Theta(n\log n)$ \\
\hline
$\displaystyle\sum_{i=1}^{n} i$ & $\Theta(n^2)$ \\
\hline
$\displaystyle\sum_{i=1}^{n} \frac1i$ & $\Theta(\log n)$ \\
\hline
$\displaystyle\sum_{i=1}^{n} \log i$ & $\Theta(n\log n)$ \\
\hline
\end{tabular}
\end{center}

\noindent In the example of \texttt{BubbleSort}, the worst case runtime was $T_{run}(n)=\frac{n(n-1)}{2}$. So you 
will often come across statements like ``\texttt{BubbleSort} has worst case time complexity $\Theta(n^2)$''. What 
this means is that the complexity of the worst case runtime as a function of $n$ is $\Theta(n^2)$.\\

\noindent A natural question which arises in comparing different algorithms for the same task is which one is faster?
For example, consider two algorithms A and B, where A has runtime $T_{run}(n)=n^2$ and B has 
runtime $T_{run}(n)=100n\log n$. To simplify our comparison, we just compare their complexities.\\


\noindent Comparing between the complexity functions $f(n)$ and $g(n)$ themselves is quite straight forward -- 
we just evaluate the limits of the ratio $f(n)/g(n)$ as $n\to\infty$. If the first 
limit is $0$, then $f(n) < g(n)$, if the limit is a non-zero constant (usually 1) then $f(n)=g(n)$, otherwise $f(n)>g(n)$.
Basis these simple rules, we can have a total order of the common complexity functions as 
\[
\Theta(1)<\Theta(\log n)<\Theta(n^\frac{1}{x})<\Theta(n)<\Theta(n\log n)<\Theta(n^2)
\]
In the above comparison, $x > 1$. Algorithms having time complexity smaller than $\Theta(n)$
are called sub-linear time algorithms.\\


\noindent Finally all the above discussion allows us to say that A: $T_{run}=n^2$ is worse than 
B: $T_{run}=100n\log n$ since A is $\Theta(n^2)$ and B is $\Theta(n\log n)$ even 
though for small n the first one may well be better, but for small n the performance doesn't really\footnote{Integer 
multiplication is an interesting case. The classic method of multiplying digit by digit 
as taught in school is $\Theta(n^2)$ time, but there are complicated methods 
which can do it in $\Theta(n\log n)$ time. But the benefit of using the complicated 
methods really doesn't show up until quite a large n because of high constant factor.}
  matter anyway\dots

\clearpage
\section{Problems}

\begin{enumerate}

\item[\textbf{1}]
A \texttt{vector} in \texttt{C++} is an array with variable size. Once the memory 
allocated to the \texttt{vector} is exceeded, it is reallocated more memory somewhere else. 
The designers of the \texttt{vector class} must choose a sequence of sizes \(s_0=0 < s_1 < s_2 < \ldots\) 
such that when the size of the \texttt{vector} becomes any $s_i$, we reallocate and assign memory of size 
\(s_{i+1}\) to the \texttt{vector}.\\


\noindent Thus each \texttt{push\_back()} operation is one operation if space is available and \(s_i+1\) otherwise. 
For the average case analysis take all sizes of vectors between $0$ and $n-1$ (both inclusive) to be equally likely.

\begin{enumerate}[label=\textbf{(\alph*)}]
\item If $s_i=i$, show that the time complexity for the average \texttt{push\_back()} is $\Theta(n)$.
\item Show that in general the time complexity for the average \texttt{push\_back()}is $\Theta \left( 1+\frac{\left(\sum\limits_{s_i<n}s_i\right)}{n} \right)$. 
\item Knowing the previous result, come up with a simple series of \(s_i\) that would 
give a time complexity of $\Theta(1)$ for the average \texttt{push\_back()}.
\end{enumerate}


\item[\textbf{2}]
Consider the algorithm

\begin{center}
\begin{minipage}{0.9\linewidth}
\begin{framed}
\begin{Verbatim}
f(n) 
{
    count = 0
    while(n>1)
    {
        if(n even)
            n = n/2
        else //n is odd
            n = n-1
        count = count + 1
    }
    return count
}
\end{Verbatim}
\vspace{-4mm}
\end{framed}
\end{minipage}
\end{center}

Here \(f(n)\) counts the number of iterations required to reduce the number \(n\) to 1. We can think of $f(n)$ as a measure 
of the time taken by the algorithm (one iteration being one operation). Now consider $x \in \{2^{k+1}-1 : k \in \mathbb{N}\}$ and $n$ be the input constrained to be between $1$ and $x$.
Find the asymptotic\footnote{We have seen this ratio being 2 in the case of \texttt{BubbleSort}.} ratio of worst case time and average case time of the algorithm with respect to n. Equivalently, find  
\[
\lim_{x \to \infty}
\frac{\max_{1 \le n \le x} f(n)}
{\left( \frac{1}{x} \sum_{n=1}^{x} f(n) \right)}
\]

\end{enumerate}